% Generated draft: requires mathematical and source review.
\documentclass[11pt]{article}
\usepackage[margin=25mm]{geometry}
\usepackage{fontspec}
\setmainfont{Latin Modern Roman}
\usepackage{amsmath,amssymb,mathtools,unicode-math}
\setmathfont{Latin Modern Math}
\usepackage{enumitem,needspace,xurl,hyperref}
\setlength{\parindent}{0pt}
\setlength{\parskip}{5pt}
\setlength{\emergencystretch}{3em}
\newcommand{\N}{\mathbb{N}}
\newcommand{\Z}{\mathbb{Z}}
\newcommand{\Q}{\mathbb{Q}}
\newcommand{\R}{\mathbb{R}}
\newcommand{\C}{\mathbb{C}}
\newcommand{\F}{\mathbb{F}}
\newcommand{\A}{\mathbb{A}}
\newcommand{\PP}{\mathbb{P}}
\newcommand{\E}{\mathbb{E}}
\newcommand{\eps}{\varepsilon}
\newcommand{\dd}{\,\mathrm{d}}
\begin{document}
\section*{UnsolvedMath ID: 30000558}
\textbf{Source problem code:} catalogue entry
\subsection{Definitions and mathematical statement}
For the intended integer \(k\), let \(M\) be a compact manifold of dimension \(k+1\). A \(C^{1+\varepsilon}\) action of \(\Z^k\) is a homomorphism
\[
 \alpha:\Z^k\longrightarrow\operatorname{Diff}^{1+\varepsilon}(M),
\]
with the usual H\"older-derivative interpretation when \(0<\varepsilon\leq1\). The source asks which such manifolds admit an action with no proportional Lyapunov exponents and with some \(a\in\Z^k\) of positive entropy. The rank range, fixed-versus-existential \(\varepsilon\), boundary convention and requirement of faithfulness are not supplied.

Lyapunov exponents require an invariant-measure or other specified regular-point setup. For example, with a common invariant ergodic probability measure \(\mu\), the multiparameter multiplicative ergodic theorem gives exponent functionals \(\chi_i\in(\R^k)^*\) and a measurable splitting into spaces \(E_i(x)\), with the corresponding growth formula
\[
 \lim_{n\to\infty}\frac1n\log\|D_x\alpha(na)v\|=\chi_i(a)
 \quad\text{for }0\ne v\in E_i(x)
\]
at regular points and for each \(a\in\Z^k\). The record does not specify such a measure, ergodicity, the scope in regular points, or whether the condition is required for one measure or every invariant measure. These are not optional data for a measure-theoretic interpretation.

For nonzero exponent functionals one possible reading of no proportionality is
\[
 \chi_i\ne c\chi_j\quad\text{for }i\ne j\text{ and all }c\in\R\setminus\{0\}.
\]
Excluding only positive proportionality is weaker. Zero exponents, repeated functionals and multiplicities also require conventions. Pairwise nonproportionality does not by itself specify general position of all Lyapunov hyperplanes, so no stronger general-position, one-dimensional-splitting or Anosov condition is added.

The entropy condition could mean \(h_\mu(\alpha(a))>0\) for the same measure used for the exponents, or \(h_{\mathrm{top}}(\alpha(a))>0\). The source does not choose; the existential element is retained without conflating those predicates. The classification must therefore first specify these conventions, then characterize the manifolds supporting the resulting actions. It is also unclear whether classification is up to homeomorphism or diffeomorphism, or permits finite covers; no toral or algebraic model is selected in advance.

The literature triage checked 2026-08-21 reports no authoritative resolution of the extracted claim located in its search. This open label and the editorial verification note are retained as dated metadata, not an independent rigidity theorem or current classification assessment.
\subsection{Short English statement}
Classify compact manifolds of dimension one more than the action rank that support the stated smooth abelian dynamics. The source omits the invariant measure defining Lyapunov exponents and the entropy convention. Nonproportional exponents are not silently strengthened to general position, simple splitting or uniform hyperbolicity.
\subsection{Sources}
\begin{itemize}
\item Local UnsolvedMath registry, permanent ID 30000558, source code catalogue entry, ``Maximal-Rank Smooth Abelian Actions,'' including original and cleaned manifold-classification question and verification metadata.
\item \emph{Geometric Group Theory, Hyperbolic Dynamics and Symplectic Geometry (2007)}, Oberwolfach Reports, pp.\ 1991--2058, DOI 10.4171/owr/2006/33: \url{https://doi.org/10.4171/owr/2006/33}. This is also the sole supplied literature reference. The citation uses 2007, while the DOI and proposed-year metadata use 2006.
\item Literature triage checked 2026-08-21; local record updated 2026-09-16. No external source has been visited or independently checked here.
\end{itemize}
\end{document}
